\documentclass[11pt,a4paper]{article}

% ─── Package Imports ────────────────────────────────────────────────────────
\usepackage[utf8]{inputenc}
\usepackage[T1]{fontenc}
\usepackage[left=1in,right=1in,top=1in,bottom=1in]{geometry}
\usepackage{graphicx}
\usepackage{amsmath}
\usepackage{amssymb}
\usepackage{cite}
\usepackage{hyperref}
\usepackage{booktabs}
\usepackage{float}
\usepackage{multirow}
\usepackage{array}
\usepackage{xcolor}
\usepackage{fancyhdr}
\usepackage{setspace}
\usepackage{caption}
\usepackage{subcaption}
\usepackage{enumitem}
\usepackage{url}

% ─── Page Style ─────────────────────────────────────────────────────────────
\pagestyle{fancy}
\fancyhf{}
\rhead{\thepage}
\lhead{YOLO11n-Ghost-Hybrid for Insulator Defect Detection}
\setstretch{1.5}

\hypersetup{
    colorlinks=true,
    linkcolor=black,
    citecolor=black,
    urlcolor=black
}

% ─── Title ───────────────────────────────────────────────────────────────────
\title{%
  \LARGE\textbf{YOLO11n-Ghost-Hybrid: An Ultra-Lightweight \\
  Knowledge-Distilled Architecture for Edge-Deployed \\
  Insulator Defect Detection}%
}

\author{%
  P~P~Satya~Karthikeya \and
  B~Karthikeya \and
  M~Karthik~Reddy \and
  P~Rohit%
}

\date{\today}

% ═══════════════════════════════════════════════════════════════════════════════
\begin{document}
\maketitle
\thispagestyle{empty}

% ─── Abstract ────────────────────────────────────────────────────────────────
\begin{abstract}
\noindent
Real-time insulator defect detection on edge hardware is fundamentally constrained by model
size, compute budget, and class imbalance between normal and damaged components. This paper
presents \textit{YOLO11n-Ghost-Hybrid-P3P4-Medium}, a domain-specialised lightweight
detector for two-class insulator defect detection (\texttt{Damaged\_1} and
\texttt{insulator}). The proposed architecture replaces the standard YOLO11n backbone and
detection head with GhostConv/C3Ghost stages and a depthwise-separable FPN-PAN head,
respectively, and removes the large-scale P5 detection branch to focus all capacity on
small and medium defect targets. Training follows a three-stage pipeline: (i) from-scratch
architecture training, (ii) teacher-student knowledge distillation (KD) from a YOLO11s
teacher, and (iii) resolution-aware fine-tuning at 768$\times$768. The final model
contains \textbf{867,664 parameters} and \textbf{5.7 GFLOPs}, achieving \textbf{96.10\%
mAP@0.5}, \textbf{66.79\% mAP@0.5:0.95}, \textbf{93.65\% precision}, and \textbf{94.18\%
recall}---within 0.44 mAP@0.5 points of a teacher that is 10.9$\times$ larger. On a
Raspberry Pi CPU with ONNX Runtime, mean model inference time is 656.0\,ms (1.4\,FPS full
pipeline), demonstrating practical field deployability. A comprehensive cumulative ablation
study identifies each design decision's individual contribution, showing that architecture
redesign accounts for the largest single gain (+4.01 mAP@0.5), while resolution-aware
fine-tuning at 768 provides the highest-ROI optimisation step (+1.59 mAP@0.5).

\vspace{0.5em}
\noindent\textbf{Keywords:} YOLO11n, GhostConv, knowledge distillation, insulator defect
detection, edge AI, lightweight object detection, Raspberry Pi deployment,
resolution-aware fine-tuning
\end{abstract}

\newpage
\tableofcontents
\newpage

% ═══════════════════════════════════════════════════════════════════════════════
\section{Introduction}
\label{sec:intro}

Power-line insulator monitoring is a high-impact computer vision problem in which missed
detections can lead to delayed maintenance and significant risks to grid reliability.
Insulators are exposed to mechanical stress, corona discharge, pollution, and UV
degradation, making periodic automated inspection critical. Defect regions are typically
small, texture-like, and visually subtle---qualities that make them difficult to detect
with standard detectors and practically impossible to localise reliably with large models
that must run on low-power embedded hardware.

Single-stage detectors in the YOLO family are attractive for this setting because of their
favourable speed-accuracy trade-off. Standard compact variants, however, leave room for
further optimisation when the deployment target is a Raspberry Pi or Jetson-class embedded
system with tight memory and compute budgets. Two challenges are particularly acute:

\begin{enumerate}[leftmargin=2em]
  \item \textbf{Class imbalance.} The defect class (\texttt{Damaged\_1}) appears at
    roughly 1:3.5 frequency relative to normal insulators. Without targeted optimisation,
    recall on the rare class saturates early.
  \item \textbf{Scale mismatch.} Generic YOLO detectors allocate parameters and
    receptive field across three scales (P3/P4/P5). For insulator datasets whose targets
    are exclusively small-to-medium, the large-scale P5 branch wastes capacity.
\end{enumerate}

This paper addresses both challenges with a domain-specialised lightweight architecture
and a staged optimisation pipeline. Starting from a vanilla YOLO11n baseline (90.20\%
mAP@0.5), we progressively apply architecture-level compression by design, teacher-guided
representation transfer, and resolution-aware fine-tuning to achieve 96.10\% mAP@0.5 with
a 3$\times$ smaller model.

\subsection*{Major Contributions}

\begin{enumerate}[leftmargin=2em]
  \item We design an ultra-light YOLO11n variant (867,664 parameters) using
    GhostConv/C3Ghost and a fully depthwise detection head, optimised for
    small/medium insulator defects.
  \item We demonstrate that removing the P5 detection branch improves task-relevant
    accuracy under strict parameter budgets.
  \item We validate that KD enables stronger downstream adaptation to
    high-resolution fine-tuning, even when direct KD gain is modest.
  \item We identify 768$\times$768 as the optimal operating resolution for this
    architecture family, delivering the largest single training-step gain.
  \item We provide a rigorous cumulative ablation study (sourced from
    \texttt{docs/ABLATION\_STUDY\_MAIN\_MODEL.md}) that quantifies each design
    decision's contribution in isolation.
  \item We demonstrate edge feasibility on Raspberry Pi via ONNX Runtime with
    reproducible latency statistics and a complete deployment workflow.
\end{enumerate}

% ═══════════════════════════════════════════════════════════════════════════════
\section{Related Work}
\label{sec:related}

\subsection{One-Stage Detection for Real-Time Applications}

YOLO-style one-stage detectors unify localisation and classification in a single forward
pass, making them suitable for latency-sensitive deployments~\cite{redmon2016yolo,
bochkovskiy2020yolov4, wang2022yolov7}. Recent generations (YOLOv8/YOLO11) improve
training stability and multi-scale detection quality, but edge-focused tuning remains
application-dependent~\cite{ultralytics2023yolov8, ultralytics2024yolo11}.

\subsection{Lightweight CNN Design}

Model families such as MobileNet and GhostNet show that depthwise-separable operations
and cheap feature generation can preserve representational power while significantly
reducing multiply-accumulate (MAC)
operations~\cite{howard2017mobilenets, han2020ghostnet}. GhostNet's key insight---that
many feature maps are redundant and can be approximated via cheap linear
operations---directly motivates our backbone design.

\subsection{Knowledge Distillation for Compact Models}

Knowledge distillation (KD) transfers soft-target structure from larger teacher models
into smaller students, often improving compact model calibration and rare-class
behaviour~\cite{hinton2015distilling, wang2022kd_survey}. In this work KD alone provides
moderate direct gains but critically enables stronger downstream resolution adaptation by
improving internal representation quality before the high-resolution fine-tuning stage.

\subsection{Insulator Defect Detection}

Automated inspection of power-line insulators using deep learning has been explored with
CNN classifiers, Faster R-CNN variants, and more recently YOLO-family detectors. The
primary challenge reported across these studies is the low base rate of defect instances
and the high visual similarity between defect types and normal insulator texture under
varying environmental conditions. Our work addresses this directly through architecture
design and training curriculum.

\subsection{Edge Deployment and Practical Constraints}

Prior edge-AI studies emphasise that theoretical complexity does not always translate to
real-time throughput without careful runtime and format optimisation (ONNX / TensorRT /
TFLite)~\cite{tan2019efficientnet, onnxruntime2024}. Deployment choices are therefore
part of model design, not an afterthought, which is why we benchmark on Raspberry Pi
hardware rather than on the training GPU.

% ═══════════════════════════════════════════════════════════════════════════════
\section{Dataset}
\label{sec:dataset}

\subsection{Composition}

The dataset covers two classes: \texttt{Damaged\_1} (the rare defect class, comprising
small cracks, burns, and surface degradation) and \texttt{insulator} (the normal class).
Table~\ref{tab:dataset} summarises the split statistics.

\begin{table}[H]
\centering
\caption{Dataset split statistics}
\label{tab:dataset}
\begin{tabular}{lrrrr}
\toprule
\textbf{Split} & \textbf{Images} &
\textbf{Damaged\_1} & \textbf{insulator} & \textbf{Total instances} \\
\midrule
Train & 1{,}155 & 562  & 1{,}962 & 2{,}524 \\
Val   &   143   &  76  &   224   &   300   \\
\bottomrule
\end{tabular}
\end{table}

The class ratio is approximately 1:3.5 (\texttt{Damaged\_1}:\texttt{insulator}), creating
a moderate imbalance that can suppress rare-defect recall if optimisation is not targeted.
Labels are stored in YOLO \texttt{txt} format; VOC XML annotations are also available for
compatibility. Dataset configuration is managed via \texttt{VOC/voc.yaml}.

\subsection{Augmentation Strategy}

Training uses moderate geometric and photometric augmentation: mosaic (0.8--1.0), mixup
(0.05--0.10), random rotation ($\pm$5--10\,$^\circ$), translation (0.1--0.2), scale
(0.3--0.9), and HSV jitter ($h$:0.015, $s$:0.5--0.7, $v$:0.3--0.4). For stable
late-stage convergence, heavy operations (e.g., large-scale mosaic) are reduced during
fine-tuning phases. Aggressive synthetic operations such as copy-paste and \texttt{flipud}
were deliberately excluded because they destabilised class balance in preliminary
experiments.

% ═══════════════════════════════════════════════════════════════════════════════
\section{Proposed Architecture}
\label{sec:arch}

\subsection{Design Goal and Baseline}

The reference baseline is vanilla YOLO11n (2.58M parameters, 90.20\% mAP@0.5). Our goal
is to obtain comparable or superior detection quality with a student model whose parameter
count is reduced by approximately 3$\times$, while keeping inference latency on CPU-only
edge hardware within practical bounds for periodic inspection workflows.

\subsection{YOLO11n-Ghost-Hybrid-P3P4-Medium}

Table~\ref{tab:design_choices} lists the three core architectural decisions and their
rationale.

\begin{table}[H]
\centering
\caption{Core design choices of YOLO11n-Ghost-Hybrid-P3P4-Medium}
\label{tab:design_choices}
\begin{tabular}{lll}
\toprule
\textbf{Design choice} & \textbf{What it replaces} & \textbf{Motivation} \\
\midrule
GhostConv backbone    & Standard Conv backbone      & $\approx$50\% fewer backbone parameters \\
DWConv-only head      & Standard Conv head          & $\sim$9$\times$ parameter reduction in head \\
P3+P4 detection only  & 3-scale (P3/P4/P5) head     & All capacity focused on small/medium defects \\
\bottomrule
\end{tabular}
\end{table}

The overall topology is:
\begin{itemize}[leftmargin=2em]
  \item \textbf{Stem}: standard Conv(32, 3$\times$3, $s$=2) for initial edge extraction.
  \item \textbf{Backbone}: GhostConv downsampling at P2/P3/P4/P5 (channel progression
    64$\to$128$\to$160$\to$160) with C3Ghost blocks ($\times$2 at P2, $\times$3 at
    P3/P4, $\times$2 at P5) and SPPF(160, $k$=5) context aggregation.
  \item \textbf{Neck}: FPN top-down path (P5$\to$P4$\to$P3) and PAN bottom-up path
    (P3$\to$P4), realised entirely with DWConv + pointwise Conv(1$\times$1) blocks and
    nearest-neighbour upsampling.
  \item \textbf{Detection head}: Detect([P3, P4], nc=2)---anchor-free, DFL regression.
    P5 is used only as context (not for prediction) and is subsequently discarded.
\end{itemize}

The resulting model has \textbf{867,664 parameters} and \textbf{5.7 GFLOPs} at
768$\times$768 input (Table~\ref{tab:param_budget}).

\begin{table}[H]
\centering
\caption{Parameter budget breakdown}
\label{tab:param_budget}
\begin{tabular}{lrr}
\toprule
\textbf{Component}                   & \textbf{Approx.\ params} & \textbf{\% of total} \\
\midrule
Stem Conv (L0)                        & $\approx$2.8K  &  0.3\% \\
GhostConv stages (L1, 3, 5, 7)       & $\approx$180K  & 20.8\% \\
C3Ghost blocks (L2, 4, 6, 8)         & $\approx$450K  & 51.9\% \\
SPPF (L9)                             & $\approx$41K   &  4.7\% \\
Head DWConv + Conv (L12--22)          & $\approx$190K  & 21.9\% \\
Detect layer                          & $\approx$3K    &  0.3\% \\
\midrule
\textbf{Total}                        & \textbf{$\approx$867K} & \textbf{100\%} \\
\bottomrule
\end{tabular}
\end{table}

\subsection{Mathematical Formulation}

\subsubsection{Ghost Convolution}

For a standard convolution with $C_{in}$ input channels, $C_{out}$ output channels,
kernel $k$, and spatial resolution $H\!\times\!W$, the MAC complexity is
\begin{equation}
  \mathrm{MAC}_{\mathrm{std}} = C_{in}\,C_{out}\,k^2\,H\,W.
\end{equation}
Ghost-style generation reduces heavy convolution usage by producing only $m = C_{out}/s$
intrinsic maps with full convolution and generating the remaining maps via cheap linear
transforms:
\begin{equation}
  \mathrm{MAC}_{\mathrm{ghost}} \approx C_{in}\,m\,k^2\,H\,W + m\,s\,d^2\,H\,W,
  \qquad m = \tfrac{C_{out}}{s},
\end{equation}
where $s$ is the ghost ratio and $d$ is the cheap-transform kernel size.

\subsubsection{Depthwise-Separable Detection Head}

Head blocks decompose convolution into a depthwise stage and a pointwise stage:
\begin{equation}
  \mathrm{Conv}_{dw+pw}
    = \mathrm{Conv}_{dw}(k\!\times\!k) + \mathrm{Conv}_{pw}(1\!\times\!1),
\end{equation}
achieving a parameter reduction of roughly $k^2 / (1 + k^2/C_{out})$ relative to a
full $k\!\times\!k$ convolution for the same channel count.

\subsubsection{Why P3/P4-Only Detection}

\begin{table}[H]
\centering
\caption{Scale analysis: why P5 is removed}
\label{tab:p5_analysis}
\begin{tabular}{lllll}
\toprule
\textbf{Scale} & \textbf{Stride} & \textbf{Recept.\ field} & \textbf{Best for} & \textbf{This dataset?} \\
\midrule
P3 &  8 & $\approx$50\,px  & Small objects        & \checkmark\ \texttt{Damaged\_1} cracks \\
P4 & 16 & $\approx$100\,px & Medium objects       & \checkmark\ Insulators \\
P5 & 32 & $\approx$200\,px & Large objects (cars) & \texttimes\ \textbf{Removed} \\
\bottomrule
\end{tabular}
\end{table}

Removing P5 eliminates approximately 15\% of head parameters, removes false positives
from large-scale anchors on tiny targets, and lets the model specialise entirely on
small/medium detection.

% ═══════════════════════════════════════════════════════════════════════════════
\section{Training Methodology}
\label{sec:training}

\subsection{Hardware and Experimental Setup}

All training is performed on a single NVIDIA RTX 3050 Laptop GPU (4\,GB VRAM) running
Fedora Linux with Python~3.11, PyTorch~2.9.1, and Ultralytics~8.3.240. Mandatory
constraints driven by the 4\,GB VRAM limit are:
\begin{itemize}[leftmargin=2em]
  \item \texttt{batch $\leq$ 8} at 640; drop to 4 at 704, 2 at 768, 1 at 896.
  \item \texttt{amp = True} (mixed-precision training throughout).
  \item \texttt{cache = "ram"} (no disk caching to avoid I/O instability).
  \item \texttt{workers = 4} (prevent RAM thrashing on 16\,GB system RAM).
\end{itemize}

\subsection{Stage 1 --- From-Scratch Architecture Training (exp\_002)}

The Ghost-Hybrid student model is trained from scratch for 300 epochs at
640$\times$640 with AdamW optimiser ($\text{lr}_0 = 0.01$, $\text{lrf} = 0.1$), full
augmentation recipe, and \texttt{batch = 8}. This produces the architecture baseline
(exp\_002: 94.21\% mAP@0.5). A prior run with auto-batch and disk caching achieved only
80.78\% mAP@0.5 from the identical architecture, confirming that training stability is
highly sensitive to batch and caching configuration.

\subsection{Stage 2 --- Teacher Training and Knowledge Distillation (exp\_004, exp\_005)}

A YOLO11s teacher (9.43M parameters, 21.3 GFLOPs) is trained for 150 epochs at
640$\times$640 (\texttt{batch = 4}, $\text{lr}_0 = 0.002$), reaching 96.54\% mAP@0.5
and establishing the performance ceiling. Knowledge distillation from teacher to student
uses:
\begin{equation}
  \mathcal{L}_{\mathrm{total}}
    = (1-\alpha)\,\mathcal{L}_{\mathrm{yolo}}
      + \alpha\,T^2\,\mathrm{KL}\!\left(
          \sigma\!\left(\tfrac{z_t}{T}\right)
          \Big\|
          \sigma\!\left(\tfrac{z_s}{T}\right)
        \right),
\end{equation}
where $z_t$, $z_s$ are teacher and student logits, $T$ is temperature, $\alpha$ is the
distillation weight, and $\mathcal{L}_{\mathrm{yolo}}$ contains box, class, and
distribution focal loss (DFL) terms. The KD run (exp\_005) uses $\alpha = 0.5$,
$T = 4.0$, \texttt{batch = 2}, producing 94.53\% mAP@0.5---a direct gain of +0.32 over
exp\_002, with importantly improved internal representations for the subsequent
fine-tuning stage.

\subsection{Stage 3 --- Resolution-Aware Fine-Tuning (exp\_009, exp\_012)}

Fine-tuning the KD checkpoint at 768$\times$768 (44\% more pixels than 640) with SGD
($\text{lr}_0 = 0.001$, \texttt{batch = 2}) for 25 epochs produces exp\_009: 96.11\%
mAP@0.5, a +1.59 gain over the KD baseline. This is the highest-ROI single optimisation
step. A subsequent head-only fine-tuning run (\texttt{freeze = 10}, 30 epochs) at the
same resolution produces exp\_012: 96.46\% mAP@0.5 (training-log peak) / 96.10\%
validated (no TTA).

\subsubsection{Negative Resolution Experiments}

Several alternative resolution strategies were evaluated as controls and found to be
inferior (Table~\ref{tab:neg_controls}), confirming that 768 is the resolution sweet spot
for this architecture.

\begin{table}[H]
\centering
\caption{Negative/control resolution experiments (reference: A2 = 94.53\%)}
\label{tab:neg_controls}
\begin{tabular}{llrr}
\toprule
\textbf{Control} & \textbf{Setting} & \textbf{mAP@0.5} & \textbf{$\Delta$ vs A2} \\
\midrule
C1 & 704 aggressive ($\text{lr}_0{=}0.005$, unfrozen) & 93.08 & $-1.44$ \\
C2 & 704 frozen (\texttt{freeze=10})                  & 93.37 & $-1.16$ \\
C3 & 704 stage-2 unfrozen continuation                & 92.85 & $-1.68$ \\
C4 & 768 multi-scale continuation                     & 95.69 & $-0.84$ vs A3 \\
C5 & 768 ultra-low LR continuation                    & 96.06 & $-0.05$ vs A3 \\
C6 & 896 resolution attempt                           & 94.37 & $-1.74$ vs A3 \\
\bottomrule
\end{tabular}
\end{table}

% ═══════════════════════════════════════════════════════════════════════════════
\section{Results and Discussion}
\label{sec:results}

\subsection{Final Validated Model Metrics}

The final single-model deployment configuration (exp\_012, no TTA) is validated on the
143-image validation split using Ultralytics \texttt{model.val()}. Results are presented
in Table~\ref{tab:final_metrics}.

\begin{table}[H]
\centering
\caption{Final validated metrics --- YOLO11n-Ghost-Hybrid-P3P4-Medium}
\label{tab:final_metrics}
\begin{tabular}{lr}
\toprule
\textbf{Metric}              & \textbf{Value}  \\
\midrule
mAP@0.5 (all classes)        & 96.10\%         \\
mAP@0.5:0.95 (all classes)   & 66.79\%         \\
Precision                    & 93.65\%         \\
Recall                       & 94.18\%         \\
\texttt{Damaged\_1} AP@0.5   & 95.05\%         \\
\texttt{insulator} AP@0.5    & 97.15\%         \\
Parameters                   & 867{,}664       \\
GFLOPs                       & 5.7             \\
Input resolution             & 768$\times$768  \\
\bottomrule
\end{tabular}
\end{table}

\subsection{Comparison with Baseline and Teacher}

Table~\ref{tab:comparison} compares the proposed student model against the vanilla
YOLO11n baseline and the YOLO11s teacher.

\begin{table}[H]
\centering
\caption{Student vs.\ baseline vs.\ teacher performance comparison}
\label{tab:comparison}
\begin{tabular}{lrrrrrr}
\toprule
\textbf{Model}           & \textbf{mAP@0.5} & \textbf{mAP@0.5:0.95} &
\textbf{Prec.} & \textbf{Recall} & \textbf{Params} & \textbf{GFLOPs} \\
\midrule
YOLO11n (baseline)       & 90.20\%  & 58.19\% & 89.70\% & 80.30\% & 2.58M  & ---   \\
\textbf{Ours (student)}  & \textbf{96.10\%} & \textbf{66.79\%} & \textbf{93.65\%} & \textbf{94.18\%} & \textbf{0.867M} & \textbf{5.7} \\
YOLO11s (teacher)        & 96.54\%  & ---     & ---     & ---     & 9.43M  & 21.3  \\
\bottomrule
\end{tabular}
\end{table}

The proposed model improves upon the vanilla baseline by \textbf{+5.90 mAP@0.5 points}
while using roughly $3\times$ fewer parameters. It approaches teacher-level performance
within \textbf{0.44 mAP@0.5 points} while being \textbf{10.9$\times$ smaller} and
consuming \textbf{3.7$\times$ fewer FLOPs}, demonstrating near-perfect knowledge transfer
efficiency.

\subsection{Per-Class Analysis}

\begin{table}[H]
\centering
\caption{Per-class detection performance of the final model}
\label{tab:per_class}
\begin{tabular}{lrrrr}
\toprule
\textbf{Class}               & \textbf{AP@0.5} & \textbf{AP@0.5:0.95} &
\textbf{Precision} & \textbf{Recall} \\
\midrule
\texttt{Damaged\_1} (defect) & 95.05\% & 64.9\% & 94.5\% & 90.8\% \\
\texttt{insulator} (normal)  & 97.15\% & 68.7\% & 92.8\% & 97.6\% \\
\bottomrule
\end{tabular}
\end{table}

Achieving $>$95\% AP@0.5 on the rare \texttt{Damaged\_1} class with only 867K parameters
is the key outcome. It demonstrates that the combination of architecture design, KD, and
resolution fine-tuning is sufficient to overcome class imbalance without requiring
specialised loss re-weighting or synthetic augmentation.

\subsection{Qualitative Observations}

Visual inspection of validation predictions indicates improved localisation of thin
crack-like patterns on damaged insulators, a reduced miss rate for rare defects after
KD\,+\,768 fine-tuning, and fewer high-confidence false negatives compared with the
baseline.

% ═══════════════════════════════════════════════════════════════════════════════
\section{Ablation Study}
\label{sec:ablation}

\textit{Source: \texttt{docs/ABLATION\_STUDY\_MAIN\_MODEL.md}. All values are
\textbf{peak per-epoch validation metrics} taken from \texttt{experiments/*/results.csv};
for every run the maximum epoch value is used.}

\subsection{Cumulative Main-Path Ablation}

Table~\ref{tab:cumulative_ablation} traces the cumulative contribution of each design
and training decision to the final model checkpoint
(\texttt{exp\_012\_head\_finetune\_768}).

\begin{table}[H]
\centering
\caption{Cumulative ablation along the main optimisation path (source: \texttt{docs/ABLATION\_STUDY\_MAIN\_MODEL.md})}
\label{tab:cumulative_ablation}
\begin{tabular}{clllrrrr}
\toprule
\textbf{Step} & \textbf{Change introduced} & \textbf{Run} &
\textbf{mAP@0.5} & \textbf{$\Delta$} &
\textbf{mAP@0.5:0.95} & \textbf{Precision} & \textbf{Recall} \\
\midrule
A0 & Vanilla baseline (YOLO11n) &
  \texttt{baseline\_yolo11n} & 90.20 & --- & 59.31 & 91.28 & 85.14 \\
A1 & Ghost-Hybrid-P3P4 architecture &
  \texttt{exp\_002\_ghost\_hybrid\_medium3} & 94.21 & \textbf{+4.01} & 68.08 & 94.26 & 92.45 \\
A2 & Knowledge distillation &
  \texttt{exp\_005\_kd\_student3} & 94.53 & +0.31 & 68.70 & 93.46 & 92.88 \\
A3 & 768 fine-tune (SGD, unfrozen) &
  \texttt{exp\_009\_finetune\_768} & 96.11 & \textbf{+1.59} & 66.02 & 93.80 & 95.36 \\
A4 & Head-only fine-tune at 768 (\texttt{freeze=10}) &
  \texttt{exp\_012\_head\_finetune\_768} & 96.46 & +0.35 & 66.87 & 95.08 & 95.16 \\
\bottomrule
\end{tabular}
\end{table}

\textbf{Net gain (A0$\to$A4):} +6.26 mAP@0.5 points (90.20 $\to$ 96.46). The largest
single gain comes from architecture redesign (A1: +4.01), followed by 768 fine-tuning
(A3: +1.59).

\subsection{Main Findings from the Ablation}

\begin{enumerate}[leftmargin=2em]
  \item \textbf{Architecture redesign (A0$\to$A1) delivers the largest foundational
    jump.} Replacing vanilla YOLO11n with Ghost-Hybrid-P3P4 yields +4.01 mAP@0.5
    while simultaneously reducing parameter count by $3\times$. This confirms that
    task-aligned architectural compression is more effective than post-hoc pruning.

  \item \textbf{KD (A1$\to$A2) gives a modest direct gain (+0.31) but enables stronger
    high-resolution transfer.} The KD step improves the quality of learned
    representations, which is critical for downstream adaptation at 768 pixels where
    the training signal is noisier due to the forced batch-size reduction.

  \item \textbf{Resolution sweet spot is 768, not 704 or 896 (A2$\to$A3: +1.59).}
    The 768 fine-tune is the highest-ROI single step. All four resolution alternatives
    tested (704 aggressive/frozen, 768 multi-scale/ultra-low-LR, 896) underperformed
    relative to the clean 768 run (Table~\ref{tab:neg_controls}), confirming that
    resolution benefit is contingent on appropriate learning-rate and batch-size scaling.

  \item \textbf{Head-only final tuning provides the last stable push (A3$\to$A4:
    +0.35).} Freezing the backbone (\texttt{freeze=10}) and fine-tuning only the
    detection head prevents catastrophic forgetting of backbone features while allowing
    the detection layers to adapt further.

  \item \textbf{Trade-off note.} Relative to A2, A3 raises mAP@0.5 and recall
    strongly but causes a small mAP@0.5:0.95 dip (68.70$\to$66.02). A4 partially
    recovers mAP@0.5:0.95 (66.87) while retaining peak mAP@0.5, indicating that
    head-only tuning improves localisation precision after the resolution shift.

  \item \textbf{Teacher reference.} The teacher peak mAP@0.5 is 96.54; the final
    student (A4) reaches 96.46---a gap of only \textbf{0.08 mAP@0.5 points} on the
    training-log peak basis.
\end{enumerate}

% ═══════════════════════════════════════════════════════════════════════════════
\section{Edge Deployment and Hardware Validation}
\label{sec:edge}

\subsection{Export Workflow}

The trained PyTorch checkpoint (\texttt{best.pt}) is exported to ONNX (opset 17) for
Raspberry Pi CPU inference and to TensorRT FP16 for Jetson-class deployment, using the
Ultralytics CLI:
\begin{verbatim}
yolo export model=experiments/exp_012_head_finetune_768/weights/best.pt \
           format=onnx imgsz=768 opset=17
\end{verbatim}
TFLite INT8 export is also supported for additional CPU efficiency. The compact 867K
parameter footprint means all three export paths require no architectural
modifications.

\subsection{Raspberry Pi Benchmark}

Inference is benchmarked using the included
\texttt{raspi/inference\_benchmark.py} script on Raspberry Pi hardware with ONNX
Runtime CPU execution at 768$\times$768 FP32. The script performs configurable warmup
runs (default: 3), multiple timed runs per image (default: 10), and reports both
model-only and full pipeline (preprocess + infer + postprocess) latencies.
Results are summarised in Table~\ref{tab:raspi}.

\begin{table}[H]
\centering
\caption{Raspberry Pi ONNX Runtime CPU inference benchmark (768$\times$768 input)}
\label{tab:raspi}
\begin{tabular}{lrr}
\toprule
\textbf{Metric} & \textbf{Model only} & \textbf{Full pipeline} \\
\midrule
Mean     & 656.0\,ms & 690.9\,ms \\
Median   & 657.7\,ms & 692.1\,ms \\
Min      & 643.9\,ms & 677.4\,ms \\
Max      & 660.7\,ms & 697.1\,ms \\
Std dev  &   4.6\,ms & ---       \\
Throughput & ---     & 1.4\,FPS  \\
\bottomrule
\end{tabular}
\end{table}

The tight standard deviation (4.6\,ms) indicates stable, predictable latency---a
critical property for safety-relevant periodic inspection workflows. At 1.4\,FPS the
model is suitable for frame-by-frame image analysis of captured drone or UAV footage,
the dominant operational mode for power-line inspection.

\subsection{Jetson Orin Nano Path}

The repository includes TensorRT FP16 export support for Jetson-class accelerators.
Expected inference time is in the range 15--30\,ms depending on precision and runtime
path, corresponding to real-time multi-camera operation. This confirms that the
proposed model is accelerator-friendly and can support high-throughput inspection
scenarios without architectural changes.

\subsection{Post-Training Techniques}

Table~\ref{tab:post_training} summarises post-training enhancement techniques explored
after the main training pipeline. None improved upon the clean single-model best.

\begin{table}[H]
\centering
\caption{Post-training technique results}
\label{tab:post_training}
\begin{tabular}{lrl}
\toprule
\textbf{Technique}         & \textbf{mAP@0.5} & \textbf{Verdict} \\
\midrule
Model Soup (weight avg.)   & 95.9\%  & Below single-model best \\
SWA (Stochastic Weight Avg.)& 96.09\% & Matched best single-model \\
Multi-scale training       & 95.69\% & Hurt performance \\
Ultra-low LR continuation  & 96.06\% & Saturated, marginal gain \\
TTA (test-time augmentation)& 96.21\% & +0.11; latency overhead unacceptable for edge \\
\bottomrule
\end{tabular}
\end{table}

% ═══════════════════════════════════════════════════════════════════════════════
\section{Conclusion}
\label{sec:conclusion}

This paper presented YOLO11n-Ghost-Hybrid-P3P4-Medium, a lightweight and
deployment-oriented YOLO11n variant for insulator defect detection. By combining
Ghost-Hybrid architecture design, teacher-student knowledge distillation, and
resolution-aware fine-tuning at 768$\times$768, we obtained a model that is compact
and accurate: \textbf{867,664 parameters, 5.7 GFLOPs, and 96.10\% mAP@0.5}. The
cumulative ablation study (Section~\ref{sec:ablation}), drawn directly from
\texttt{docs/ABLATION\_STUDY\_MAIN\_MODEL.md}, demonstrates that architecture-aware
optimisation and a staged training curriculum together outperform brute-force scaling
under strict edge constraints. Raspberry Pi benchmarks confirm practical deployability
for periodic drone-based inspection, while Jetson-oriented export paths support
higher-throughput real-time use cases.

\subsection*{Future Work}

\begin{enumerate}[leftmargin=2em]
  \item Collect and report measured Jetson FP32/FP16/INT8 latency and memory traces
    in a standardised benchmark table.
  \item Evaluate post-training quantisation and quantisation-aware training (QAT) to
    improve CPU throughput on Raspberry Pi beyond 1.4\,FPS.
  \item Expand the defect taxonomy and evaluate domain generalisation across weather
    and time-of-day conditions.
  \item Integrate uncertainty estimation for risk-aware maintenance prioritisation.
  \item Investigate structured pruning as an alternative path to further parameter
    reduction while preserving the current mAP@0.5.
\end{enumerate}

% ═══════════════════════════════════════════════════════════════════════════════
\section*{Reproducibility}

All training scripts, configuration files, and experiment CSV logs are included in the
repository. Key resources are listed in Table~\ref{tab:repro}.

\begin{table}[H]
\centering
\caption{Reproducibility resources}
\label{tab:repro}
\begin{tabular}{ll}
\toprule
\textbf{Resource}              & \textbf{Path} \\
\midrule
Architecture YAML              & \texttt{models/yolo11n-ghost-hybrid-p3p4-medium.yaml} \\
Dataset config                 & \texttt{VOC/voc.yaml} \\
Training scripts               & \texttt{scripts/train\_*.py} \\
Baseline log                   & \texttt{experiments/baseline\_yolo11n/results.csv} \\
Experiment logs                & \texttt{experiments/exp\_*/results.csv} \\
Ablation source                & \texttt{docs/ABLATION\_STUDY\_MAIN\_MODEL.md} \\
Architecture notes             & \texttt{docs/ARCHITECTURE.md} \\
Experiment journey             & \texttt{docs/EXPERIMENT\_JOURNEY.md}, \texttt{EXPERIMENT\_JOURNEY\_V2.md} \\
Benchmark report               & \texttt{BENCHMARK.md} \\
Raspberry Pi benchmark script  & \texttt{raspi/inference\_benchmark.py} \\
\bottomrule
\end{tabular}
\end{table}

% ═══════════════════════════════════════════════════════════════════════════════
\begin{thebibliography}{99}

\bibitem{redmon2016yolo}
J. Redmon, S. Divvala, R. Girshick, and A. Farhadi,
``You Only Look Once: Unified, Real-Time Object Detection,''
in \textit{Proc.\ IEEE Conf.\ Computer Vision and Pattern Recognition (CVPR)}, 2016.

\bibitem{bochkovskiy2020yolov4}
A. Bochkovskiy, C.-Y. Wang, and H.-Y. M. Liao,
``YOLOv4: Optimal Speed and Accuracy of Object Detection,''
\textit{arXiv preprint arXiv:2004.10934}, 2020.

\bibitem{wang2022yolov7}
C.-Y. Wang, A. Bochkovskiy, and H.-Y. M. Liao,
``YOLOv7: Trainable Bag-of-Freebies Sets New State-of-the-Art for Real-Time Object
Detectors,''
\textit{arXiv preprint arXiv:2207.02696}, 2022.

\bibitem{ultralytics2023yolov8}
Ultralytics,
``YOLOv8 Documentation and Release Notes,''
\url{https://docs.ultralytics.com}, 2023--2026.

\bibitem{ultralytics2024yolo11}
Ultralytics,
``YOLO11 Documentation and Release Notes,''
\url{https://docs.ultralytics.com}, 2024--2026.

\bibitem{howard2017mobilenets}
A. G. Howard, M. Zhu, B. Chen, D. Kalenichenko, W. Wang, T. Weyand, M. Andreetto,
and H. Adam,
``MobileNets: Efficient Convolutional Neural Networks for Mobile Vision Applications,''
\textit{arXiv preprint arXiv:1704.04861}, 2017.

\bibitem{han2020ghostnet}
K. Han, Y. Wang, Q. Tian, J. Guo, C. Xu, and C. Xu,
``GhostNet: More Features from Cheap Operations,''
in \textit{Proc.\ IEEE Conf.\ Computer Vision and Pattern Recognition (CVPR)}, 2020.

\bibitem{hinton2015distilling}
G. Hinton, O. Vinyals, and J. Dean,
``Distilling the Knowledge in a Neural Network,''
\textit{arXiv preprint arXiv:1503.02531}, 2015.

\bibitem{wang2022kd_survey}
Y. Wang, W. Zhou, T. Jiang, X. Bai, and Y. Xu,
``Mimicking the Oracle: An Initial Phase Decorrelation Approach for Class Incremental
Learning,''
\textit{arXiv preprint arXiv:2112.04731}, 2022.

\bibitem{tan2019efficientnet}
M. Tan and Q. V. Le,
``EfficientNet: Rethinking Model Scaling for Convolutional Neural Networks,''
in \textit{Proc.\ Int.\ Conf.\ Machine Learning (ICML)}, 2019.

\bibitem{onnxruntime2024}
Microsoft,
``ONNX Runtime: Performance Tuning and Execution Providers,''
\url{https://onnxruntime.ai/docs/performance}, 2024--2026.

\end{thebibliography}

\end{document}
